
\documentclass[11pt,a4paper]{article}

\usepackage[margin=0.8in]{geometry}
\usepackage[T1]{fontenc}
\usepackage{lmodern}
\usepackage{xcolor}
\usepackage{hyperref}
\usepackage{enumitem}
\usepackage{titlesec}
\usepackage{tabularx}
\usepackage{fancyhdr}
\usepackage{tcolorbox}
\usepackage{amssymb}

\definecolor{darkblue}{RGB}{30,65,110}
\definecolor{lightblue}{RGB}{235,243,252}
\definecolor{lightgray}{RGB}{245,245,245}
\definecolor{darkgreen}{RGB}{35,105,70}

\hypersetup{colorlinks=true,linkcolor=darkblue,urlcolor=darkblue}

\titleformat{\section}{\Large\bfseries\color{darkblue}}{\thesection}{0.6em}{}
\titleformat{\subsection}{\large\bfseries\color{darkblue}}{\thesubsection}{0.5em}{}

\pagestyle{fancy}
\fancyhf{}
\lhead{Unit 1 -- UML}
\rhead{Course Tracker}
\cfoot{\thepage}

\newtcolorbox{conceptbox}{
    colback=lightblue,colframe=darkblue,boxrule=0.7pt,arc=2mm,
    left=2mm,right=2mm,top=1.5mm,bottom=1.5mm
}
\newtcolorbox{exambox}{
    colback=lightgray,colframe=black!40,boxrule=0.5pt,arc=2mm,
    left=2mm,right=2mm,top=1.5mm,bottom=1.5mm
}

\newcommand{\checkitem}{$\square$}

\title{\vspace{-1cm}\textbf{Unit 1 -- UML}\\
\large Complete Topic \& Subtopic Tracker}
\author{Course Notes}
\date{}

\begin{document}
\maketitle

\begin{conceptbox}
\textbf{Syllabus scope:} UML, with emphasis on the two diagram types listed
in the course syllabus: \textbf{Class Diagrams} and \textbf{Sequence Diagrams}.
\end{conceptbox}

\tableofcontents
\newpage

% =========================================================
\section{UML Fundamentals}

\subsection{Study Tracker}

\begin{itemize}[label=\checkitem]
    \item Know what UML stands for
    \item Understand why UML is used
    \item Understand structural vs interaction modeling
    \item Know the purpose of class diagrams
    \item Know the purpose of sequence diagrams
\end{itemize}

\subsection{What is UML?}

UML stands for \textbf{Unified Modeling Language}.

It is a standardized visual language used to represent and communicate
software designs.

Instead of immediately writing implementation code, a designer can first
represent the system visually.

\subsection{Why Use UML?}

UML can help developers:

\begin{itemize}
    \item understand a system before implementation,
    \item communicate designs,
    \item identify relationships between components,
    \item visualize object interactions,
    \item document a design.
\end{itemize}

\begin{conceptbox}
\textbf{Mental model:}

Class Diagram = \textbf{What the system is made of}

Sequence Diagram = \textbf{How the objects interact}
\end{conceptbox}

% =========================================================
\section{Class Diagrams}

\subsection{Study Tracker}

\begin{itemize}[label=\checkitem]
    \item Understand the purpose of a class diagram
    \item Identify classes
    \item Identify attributes
    \item Identify methods/operations
    \item Understand visibility
    \item Understand relationships
    \item Understand multiplicity
    \item Read a complete class diagram
    \item Create a basic class diagram from requirements
\end{itemize}

\subsection{Purpose of a Class Diagram}

A class diagram represents the static structure of a system.

It can show:

\begin{itemize}
    \item classes,
    \item attributes,
    \item operations/methods,
    \item relationships,
    \item visibility,
    \item multiplicity.
\end{itemize}

\subsection{Basic Class Structure}

A UML class is commonly represented using three sections:

\begin{verbatim}
+-------------------------+
|        Student          |
+-------------------------+
| - name : String         |
| - rollNo : int          |
+-------------------------+
| + register() : void     |
| + getName() : String    |
+-------------------------+
\end{verbatim}

The three sections represent:

\begin{enumerate}
    \item Class name
    \item Attributes
    \item Operations/methods
\end{enumerate}

\subsection{Class Name}

The top section contains the class name.

Examples:

\begin{verbatim}
Student
Car
Payment
ParkingSpot
BankAccount
\end{verbatim}

\subsection{Attributes}

Attributes represent the data/state associated with the class.

Example:

\begin{verbatim}
- name : String
- balance : double
- accountNumber : int
\end{verbatim}

A typical UML attribute can be represented as:

\[
\text{visibility name : type}
\]

\subsection{Methods / Operations}

Operations represent behavior provided by the class.

Example:

\begin{verbatim}
+ deposit(amount : double) : void
+ withdraw(amount : double) : bool
+ getBalance() : double
\end{verbatim}

A typical UML operation can be represented as:

\[
\text{visibility method(parameters) : returnType}
\]

% =========================================================
\section{Visibility in Class Diagrams}

\subsection{Study Tracker}

\begin{itemize}[label=\checkitem]
    \item Public visibility
    \item Private visibility
    \item Protected visibility
    \item Understand visibility symbols
    \item Interpret visibility in a diagram
\end{itemize}

\subsection{Common Visibility Symbols}

\begin{tabularx}{\textwidth}{|p{0.2\textwidth}|p{0.2\textwidth}|X|}
\hline
\textbf{Symbol} & \textbf{Visibility} & \textbf{Meaning}\\
\hline
+ & Public & Accessible through the public interface.\\
\hline
- & Private & Accessible only within the class.\\
\hline
\# & Protected & Accessible within the class and appropriate derived classes.\\
\hline
\end{tabularx}

Example:

\begin{verbatim}
+ deposit()
- balance
# validateTransaction()
\end{verbatim}

% =========================================================
\section{Relationships Between Classes}

\subsection{Study Tracker}

\begin{itemize}[label=\checkitem]
    \item Association
    \item Aggregation
    \item Composition
    \item Generalization / inheritance
    \item Understand relationship direction
    \item Understand multiplicity
\end{itemize}

\subsection{Association}

Association represents a relationship between classes.

Example:

\[
Customer \longrightarrow Order
\]

A customer can be associated with orders.

The relationship can be named to clarify its meaning.

\subsection{Aggregation}

Aggregation represents a whole-part relationship where the parts can exist
independently of the whole.

Conceptually:

\[
Department \longrightarrow Employee
\]

An employee can exist even if a particular department object is removed.

Aggregation is commonly represented with an \textbf{empty diamond} at the
whole side.

\subsection{Composition}

Composition represents a stronger whole-part relationship.

The part is strongly owned by the whole.

Example:

\[
House \longrightarrow Room
\]

In a conceptual composition relationship, the lifecycle of the part is
strongly tied to the whole.

Composition is commonly represented with a \textbf{filled diamond} at the
whole side.

\subsection{Aggregation vs Composition}

\begin{tabularx}{\textwidth}{|p{0.25\textwidth}|X|X|}
\hline
 & \textbf{Aggregation} & \textbf{Composition}\\
\hline
Relationship & Whole-part & Strong whole-part\\
\hline
Ownership & Weaker & Stronger\\
\hline
Part lifetime & Can exist independently & Strongly tied to whole\\
\hline
UML notation & Empty diamond & Filled diamond\\
\hline
\end{tabularx}

\subsection{Generalization / Inheritance}

Generalization represents an inheritance relationship.

Example:

\[
Vehicle
\rightarrow
\begin{cases}
Car\\
Bike
\end{cases}
\]

It represents an ``is-a'' relationship.

\begin{exambox}
\textbf{Important:} Do not confuse inheritance with composition.\\
Inheritance = ``is-a''.\\
Composition = ``has-a'' / strong whole-part relationship.
\end{exambox}

% =========================================================
\section{Multiplicity}

\subsection{Study Tracker}

\begin{itemize}[label=\checkitem]
    \item Understand multiplicity
    \item Know 1
    \item Know 0..1
    \item Know *
    \item Know 1..*
    \item Interpret one-to-one relationships
    \item Interpret one-to-many relationships
    \item Interpret many-to-many relationships
\end{itemize}

\subsection{What is Multiplicity?}

Multiplicity describes how many instances of one class can participate in a
relationship with another class.

Common notation:

\begin{tabularx}{\textwidth}{|p{0.2\textwidth}|X|}
\hline
\textbf{Notation} & \textbf{Meaning}\\
\hline
1 & Exactly one\\
\hline
0..1 & Zero or one\\
\hline
* & Zero or more / many\\
\hline
1..* & One or more\\
\hline
\end{tabularx}

\subsection{One-to-One}

\[
Person \quad 1 \longleftrightarrow 1 \quad Passport
\]

One person is associated with one passport in the modeled system.

\subsection{One-to-Many}

\[
Customer \quad 1 \longrightarrow * \quad Order
\]

One customer can have many orders.

\subsection{Many-to-Many}

\[
Student \quad * \longleftrightarrow * \quad Course
\]

Many students can enroll in many courses.

\begin{exambox}
When reading multiplicity, always ask:
\textit{``How many objects on this side can be associated with one object on
the other side?''}
\end{exambox}

% =========================================================
\section{Creating a Class Diagram from Requirements}

\subsection{Study Tracker}

\begin{itemize}[label=\checkitem]
    \item Extract nouns from requirements
    \item Identify candidate classes
    \item Identify attributes
    \item Identify behaviors
    \item Identify relationships
    \item Determine multiplicity
    \item Determine appropriate relationship type
    \item Review responsibilities
\end{itemize}

\subsection{Step 1: Identify Candidate Classes}

Suppose the requirement says:

\begin{quote}
A customer can book multiple movie tickets for a show.
\end{quote}

Candidate classes include:

\begin{verbatim}
Customer
Movie
Show
Ticket
\end{verbatim}

\subsection{Step 2: Identify Attributes}

For example:

\begin{verbatim}
Customer
- customerId
- name

Movie
- movieId
- title
\end{verbatim}

\subsection{Step 3: Identify Behavior}

Possible methods:

\begin{verbatim}
Customer
+ bookTicket()

Show
+ checkAvailability()

Ticket
+ cancel()
\end{verbatim}

\subsection{Step 4: Identify Relationships}

For example:

\[
Customer \longrightarrow Booking
\]

\[
Show \longrightarrow Booking
\]

\subsection{Step 5: Determine Multiplicity}

Ask:

\begin{center}
How many bookings can one customer make?
\end{center}

and:

\begin{center}
How many customers can be associated with one booking?
\end{center}

Then represent the appropriate multiplicities.

% =========================================================
\section{Sequence Diagrams}

\subsection{Study Tracker}

\begin{itemize}[label=\checkitem]
    \item Understand the purpose of sequence diagrams
    \item Identify actors/objects
    \item Understand lifelines
    \item Understand messages
    \item Understand message order
    \item Understand return messages
    \item Create a sequence diagram from a scenario
\end{itemize}

\subsection{Purpose}

A sequence diagram represents interactions between objects in chronological
order.

It answers:

\[
\boxed{\text{Who communicates with whom, and in what order?}}
\]

\subsection{Basic Structure}

Conceptually:

\begin{verbatim}
Customer        ATM          Bank
   |             |             |
   | withdraw()  |             |
   |------------>|             |
   |             | check()     |
   |             |------------>|
   |             |<------------|
   |<------------|             |
\end{verbatim}

Time generally moves from top to bottom.

\subsection{Actors and Objects}

A sequence diagram can contain:

\begin{itemize}
    \item actors,
    \item system objects,
    \item services,
    \item external systems.
\end{itemize}

An actor is an external participant interacting with the system.

\subsection{Lifeline}

A lifeline represents the existence/participation of an actor or object during
the interaction.

It is generally drawn as a vertical dashed line below the participant.

\subsection{Messages}

A message represents communication from one participant to another.

Examples:

\begin{verbatim}
login()
validatePayment()
createBooking()
sendNotification()
\end{verbatim}

The vertical order of messages represents their sequence.

\subsection{Return Messages}

A participant can return information to the sender.

Example:

\begin{verbatim}
Customer       PaymentService
   |                  |
   | pay()            |
   |----------------->|
   |                  |
   |<-----------------|
   | paymentSuccess   |
\end{verbatim}

% =========================================================
\section{Building a Sequence Diagram}

\subsection{Study Tracker}

\begin{itemize}[label=\checkitem]
    \item Identify participants
    \item Identify the starting actor
    \item Identify the first action
    \item Follow the complete workflow
    \item Convert actions into messages
    \item Place messages in chronological order
    \item Add responses/returns where relevant
\end{itemize}

\subsection{Example: Movie Booking}

A simplified workflow:

\[
Customer
\rightarrow
BookingSystem
\rightarrow
Show
\rightarrow
Payment
\]

Possible sequence:

\begin{enumerate}
    \item Customer searches for a movie.
    \item Booking system returns available shows.
    \item Customer selects a show.
    \item Booking system checks seat availability.
    \item Customer selects a seat.
    \item Customer makes payment.
    \item Payment service confirms payment.
    \item Booking system creates the booking.
\end{enumerate}

A sequence diagram represents this workflow through ordered messages.

% =========================================================
\section{Class Diagram vs Sequence Diagram}

\subsection{Study Tracker}

\begin{itemize}[label=\checkitem]
    \item Know what a class diagram represents
    \item Know what a sequence diagram represents
    \item Explain the difference in an exam
    \item Know when each diagram is useful
\end{itemize}

\begin{tabularx}{\textwidth}{|p{0.25\textwidth}|X|X|}
\hline
 & \textbf{Class Diagram} & \textbf{Sequence Diagram}\\
\hline
Focus & Static structure & Dynamic interaction\\
\hline
Shows & Classes, attributes, methods, relationships & Objects/actors and messages\\
\hline
Main question & ``What exists?'' & ``What happens and in what order?''\\
\hline
Time & Not the main focus & Central to the diagram\\
\hline
Useful for & Designing class structure & Modeling workflows/interactions\\
\hline
\end{tabularx}

\begin{conceptbox}
\textbf{Memory trick:}

\[
\boxed{\text{Class Diagram = Structure}}
\]

\[
\boxed{\text{Sequence Diagram = Interaction}}
\]
\end{conceptbox}

% =========================================================
\section{Common UML Design Mistakes}

\subsection{Study Tracker}

\begin{itemize}[label=\checkitem]
    \item Avoid classes with unclear responsibilities
    \item Avoid incorrect relationships
    \item Check multiplicity
    \item Distinguish aggregation and composition
    \item Distinguish inheritance and association
    \item Ensure sequence messages match the described workflow
\end{itemize}

\subsection{Wrong Relationship Type}

Do not automatically use inheritance just because two classes are related.

Ask:

\begin{center}
\textit{``Is this genuinely an is-a relationship?''}
\end{center}

If not, association or composition may be more appropriate.

\subsection{Missing Multiplicity}

A relationship without multiplicity can leave an important part of the design
ambiguous.

Always ask how many objects can participate in the relationship.

\subsection{Overloaded Class}

If one class contains many unrelated responsibilities, reconsider the class
boundaries.

This connects UML design with clean coding and SOLID.

% =========================================================
\section{Exam-Oriented Questions}

\begin{enumerate}
    \item What is UML?
    \item Why is UML used in software design?
    \item What is a class diagram?
    \item What are the three main sections of a UML class representation?
    \item What do +, -, and \# represent?
    \item What is an association?
    \item What is aggregation?
    \item What is composition?
    \item Differentiate aggregation and composition.
    \item What is generalization?
    \item What is multiplicity?
    \item Explain 1:1, 1:N, and M:N relationships.
    \item What is a sequence diagram?
    \item What is a lifeline?
    \item What is a message in a sequence diagram?
    \item How does a sequence diagram represent time/order?
    \item Differentiate a class diagram and a sequence diagram.
    \item Create a class diagram for a simple system.
    \item Create a sequence diagram for a given workflow.
\end{enumerate}

% =========================================================
\section{Final Revision Checklist}

\subsection*{UML Basics}

\begin{itemize}[label=\checkitem]
    \item I can define UML.
    \item I know why UML is used.
    \item I understand structural vs interaction modeling.
\end{itemize}

\subsection*{Class Diagrams}

\begin{itemize}[label=\checkitem]
    \item I can draw a basic class.
    \item I can identify attributes.
    \item I can identify methods.
    \item I understand public/private/protected notation.
    \item I understand association.
    \item I understand aggregation.
    \item I understand composition.
    \item I understand inheritance/generalization.
    \item I understand multiplicity.
\end{itemize}

\subsection*{Sequence Diagrams}

\begin{itemize}[label=\checkitem]
    \item I can identify actors and objects.
    \item I understand lifelines.
    \item I understand messages.
    \item I understand return messages.
    \item I can represent a workflow in chronological order.
\end{itemize}

\subsection*{Application}

\begin{itemize}[label=\checkitem]
    \item I can convert requirements into candidate classes.
    \item I can identify class relationships.
    \item I can determine multiplicity.
    \item I can create a class diagram.
    \item I can create a sequence diagram.
    \item I can explain the difference between the two.
\end{itemize}

\begin{conceptbox}
\textbf{Final mental model:}

\[
\boxed{
\text{Requirements}
\rightarrow
\text{Classes + Relationships}
\rightarrow
\text{Class Diagram}
}
\]

\[
\boxed{
\text{Scenario / Workflow}
\rightarrow
\text{Objects + Messages}
\rightarrow
\text{Sequence Diagram}
}
\]
\end{conceptbox}

\end{document}
